\subsection{Comparison of Guarantees and Vocabulary Restrictions}
\label{sec:generator-guarantee-comparison}

All families retain an executable mechanism, but each construction selects a
different subset of the vocabulary in Section~\ref{sec:generator-vocabulary}.
The original sampler uses unrestricted sums and products. Bounded-domain
contraction retains these operators on a certified invariant box with bounded
innovations; its current constructor is univariate. The Shao--Wu factor mode
saturates lag leaves inside nonlinear products and checks global sensitivity,
while the product-output mode bounds the complete product and checks the
separately derived moment envelope. Standalone linear terms remain available
in both modes. Heaps generates linear multivariate VAR equations through a
stationary parameterization and initializes their lag histories jointly.

Table~\ref{tab:four-generator-strategies} provides the cross-method summary of
these guarantees and restrictions. The original, bounded-domain, and
Shao--Wu policies use directly sampled uniform histories, while Heaps uses a
stationary Gaussian history. A verification record must therefore be read
together with its initialization and innovation laws. The numerical
comparison in Section~\ref{sec:implemented-generator-results} reports the
observed explosions, tail behaviour, and generation costs for the explicitly
recorded policies.

The executed equation after transformation is the reference mechanism.
Its AST, bindings, sampling configuration, seeds, innovation law,
initialization policy, and verification record make the process auditable.
Several trajectories can share this record. Section~\ref{sec:ground-truth-generation}
will use it to define structural and contribution-based explanation targets
over the forecast horizon.
